\documentclass[11pt,a4paper,fontset=fandol]{ctexart}

\usepackage[a4paper,top=2.25cm,bottom=2.45cm,left=2.25cm,right=2.25cm,headheight=18pt,headsep=0.55cm,footskip=24pt]{geometry}
\usepackage{amsmath,amssymb,amsthm}
\usepackage{fancyhdr}
\usepackage{lastpage}
\usepackage{enumitem}
\usepackage{titlesec}
\usepackage{xcolor}
\usepackage{hyperref}
\usepackage{bookmark}
\usepackage[most]{tcolorbox}
\usepackage{setspace}
\usepackage{array}

\hypersetup{
  colorlinks=true,
  linkcolor=blue!50!black,
  urlcolor=blue!50!black
}

\setstretch{1.13}
\setlength{\parindent}{2em}
\setlength{\parskip}{0.25em}
\allowdisplaybreaks
\setlist[enumerate]{itemsep=0.45em, topsep=0.35em, leftmargin=2.4em}
\setlist[itemize]{itemsep=0.25em, topsep=0.25em, leftmargin=1.9em}

\definecolor{mainblue}{RGB}{36,83,140}
\definecolor{lightblue}{RGB}{241,246,252}
\definecolor{lightgray}{RGB}{246,246,246}
\definecolor{rulegray}{RGB}{110,118,128}

\titleformat{\section}
  {\Large\bfseries\color{mainblue}}
  {\thesection.}
  {0.45em}
  {}
  [\vspace{0.2em}\color{rulegray}\titlerule]
\titlespacing*{\section}{0pt}{1.1em}{0.7em}

\pagestyle{fancy}
\fancyhf{}
\fancyhead[L]{\small 抽屉原理与极值思想周测 B 卷}
\fancyhead[C]{\small 组合专题：综合模型与经典构造}
\fancyhead[R]{\small 刘文博}
\fancyfoot[C]{\small 第 \thepage 页 / 共 \pageref{LastPage} 页}
\renewcommand{\headrulewidth}{0.55pt}
\renewcommand{\footrulewidth}{0.35pt}

\newtcolorbox{infobox}[1]{
  breakable,
  colback=lightblue,
  colframe=mainblue,
  boxrule=0.7pt,
  arc=1mm,
  title=\textbf{#1},
  fonttitle=\bfseries
}

\newenvironment{solution}{\par\noindent\textbf{参考解答：}}{\hfill$\square$\par}
\newcommand{\answerspace}{\vspace{2.3cm}}
\newcommand{\biganswerspace}{\vspace{3.2cm}}

\begin{document}

\thispagestyle{empty}
\begin{center}
{\fontsize{22}{28}\selectfont\bfseries 抽屉原理与极值思想周测 B 卷\par}
\vspace{0.4cm}
{\large 适合小学高年级至初中低年级数学竞赛训练\par}
\end{center}

\begin{infobox}{考试说明}
\begin{itemize}
  \item 测试时间：75 分钟
  \item 满分：100 分
  \item B 卷更强调抽屉的构造方式和极值法的严密表达，特别注意“为什么这样分组”和“上界怎样达到”。
\end{itemize}
\end{infobox}

\section{试题}

\begin{enumerate}
  \item （12 分）从 $1$ 到 $40$ 中任取多少个数，才能保证其中一定有两个数的差是 $6$ 的倍数？
  \answerspace

  \item （12 分）从 $1$ 到 $2n$ 中任取 $n+1$ 个数，证明其中一定有两个连续整数。
  \answerspace

  \item （12 分）在边长为 $3$ 的正方形内任取 $10$ 个点，证明其中一定有两个点之间的距离不超过 $\sqrt{2}$。
  \biganswerspace

  \item （12 分）从 $1$ 到 $9$ 中任取 $5$ 个互不相同的整数，证明其中一定有两个数的差不超过 $2$。
  \answerspace

  \item （12 分）从 $1$ 到 $20$ 中最多能选多少个数，使得任意两个所选出的数的差都至少为 $3$？
  \biganswerspace

  \item （12 分）从 $1$ 到 $20$ 中最多能选多少个数，使得任意两个所选出的数都没有一个是另一个的 $2$ 倍？
  \biganswerspace

  \item （14 分）从 $1$ 到 $2n$ 中任取 $n+1$ 个数，证明其中一定有两个数，一个是另一个的倍数。
  \biganswerspace

  \item （14 分）在长度为 $1$ 的线段上任取 $13$ 个点，证明其中一定有两个点之间的距离不超过 $\dfrac{1}{12}$。
  \biganswerspace
\end{enumerate}

\newpage
\section{详细解答}

\begin{enumerate}
  \item \begin{solution}
  如果两个数的差是 $6$ 的倍数，那么它们除以 $6$ 的余数相同。

  因此只需把数按除以 $6$ 的余数分类。余数共有
  \[
  0,1,2,3,4,5
  \]
  共 $6$ 类。

  若想避免出现差为 $6$ 的倍数的两个数，那么每一类里最多只能取 $1$ 个数，所以最多可以安全地取 $6$ 个数。

  再多取 $1$ 个，也就是取
  \[
  7
  \]
  个数时，就一定有两个数落在同一个余数类中，因此它们的差是 $6$ 的倍数。

  所以最少取 $7$ 个数。
  \end{solution}

  \item \begin{solution}
  把 $1$ 到 $2n$ 分成 $n$ 组：
  \[
  \{1,2\},\ \{3,4\},\ \ldots,\ \{2n-1,2n\}.
  \]

  每一组中正好是一对连续整数。

  现在任取 $n+1$ 个数，而抽屉只有 $n$ 个。由抽屉原理，至少有一组里取到了两个数。

  这一组中的两个数正好是连续整数，因此命题得证。
  \end{solution}

  \item \begin{solution}
  把边长为 $3$ 的正方形分成 $9$ 个边长为 $1$ 的小正方形。

  这 $9$ 个小正方形可以看成 $9$ 个抽屉，而 $10$ 个点是 $10$ 个物体。

  由抽屉原理，至少有一个小正方形里有至少两个点。

  同一个边长为 $1$ 的小正方形中，任意两点之间的距离都不超过它的对角线长
  \[
  \sqrt{1^2+1^2}=\sqrt{2}.
  \]

  所以一定有两个点之间的距离不超过 $\sqrt{2}$。
  \end{solution}

  \item \begin{solution}
  设这 $5$ 个数从小到大排列为
  \[
  a_1<a_2<a_3<a_4<a_5.
  \]

  如果任意两个数的差都大于 $2$，那么特别地，相邻两个数的差都至少为 $3$：
  \[
  a_2-a_1\ge 3,\quad a_3-a_2\ge 3,\quad a_4-a_3\ge 3,\quad a_5-a_4\ge 3.
  \]

  相加得到
  \[
  a_5-a_1\ge 12.
  \]

  但所有数都在 $1$ 到 $9$ 之间，所以
  \[
  a_5-a_1\le 8.
  \]

  矛盾。

  因此原假设不成立，必定有两个数的差不超过 $2$。
  \end{solution}

  \item \begin{solution}
  设选出的数从小到大排列为
  \[
  a_1<a_2<\cdots<a_k.
  \]

  因为任意两个所选数的差都至少为 $3$，所以相邻两个数的差也至少为 $3$：
  \[
  a_2-a_1\ge 3,\quad a_3-a_2\ge 3,\quad \ldots,\quad a_k-a_{k-1}\ge 3.
  \]

  相加得到
  \[
  a_k-a_1\ge 3(k-1).
  \]

  又因为 $a_1\ge 1,\ a_k\le 20$，所以
  \[
  20-1\ge a_k-a_1\ge 3(k-1).
  \]

  即
  \[
  19\ge 3(k-1),
  \]
  从而
  \[
  k\le 7.
  \]

  这个上界可以达到，例如选
  \[
  1,4,7,10,13,16,19.
  \]

  因此最多能选 $7$ 个数。
  \end{solution}

  \item \begin{solution}
  把 $1$ 到 $20$ 中的数按“奇数部分”分成若干条链：
  \[
  \{1,2,4,8,16\},\quad \{3,6,12\},\quad \{5,10,20\},\quad \{7,14\},\quad \{9,18\},
  \]
  以及
  \[
  \{11\},\ \{13\},\ \{15\},\ \{17\},\ \{19\}.
  \]

  在同一条链中，相邻两个数恰好互为 $2$ 倍。因此，为了避免出现“一个是另一个的 $2$ 倍”，在每条链中不能同时选相邻两个数。

  于是：
  \begin{itemize}
    \item 链 $\{1,2,4,8,16\}$ 最多选 $3$ 个；
    \item 链 $\{3,6,12\}$ 和 $\{5,10,20\}$ 各最多选 $2$ 个；
    \item 链 $\{7,14\}$ 和 $\{9,18\}$ 各最多选 $1$ 个；
    \item 五条只含一个数的链各最多选 $1$ 个。
  \end{itemize}

  所以总数最多为
  \[
  3+2+2+1+1+1+1+1+1+1=14.
  \]

  这个上界可以达到，例如选
  \[
  \{1,3,4,5,7,9,11,12,13,15,16,17,19,20\}.
  \]
  在每条链中都采用“隔一个取一个”的方法，因此不会出现恰好 $2$ 倍的关系。

  因此最多能选 $14$ 个数。
  \end{solution}

  \item \begin{solution}
  对于 $1$ 到 $2n$ 中的每一个整数 $m$，都可以唯一写成
  \[
  m=(\text{奇数})\times 2^k
  \]
  的形式。

  把奇数部分相同的数放进同一个抽屉中。由于不超过 $2n$ 的奇数共有
  \[
  1,3,5,\ldots,2n-1
  \]
  共 $n$ 个，所以抽屉总数是 $n$ 个。

  现在任取 $n+1$ 个数，由抽屉原理，至少有两个数落在同一个抽屉里。

  设这两个数为
  \[
  a=u\cdot 2^r,
  \qquad
  b=u\cdot 2^s,
  \]
  其中 $u$ 是同一个奇数，且不妨设 $r<s$。

  那么
  \[
  b=u\cdot 2^s=(u\cdot 2^r)\cdot 2^{s-r}=a\cdot 2^{s-r}.
  \]

  所以 $b$ 是 $a$ 的倍数。

  因此一定有两个数，一个是另一个的倍数。
  \end{solution}

  \item \begin{solution}
  把长度为 $1$ 的线段平均分成 $12$ 段，每一小段长度都是
  \[
  \frac{1}{12}.
  \]

  这 $12$ 段可以看成 $12$ 个抽屉，$13$ 个点可以看成 $13$ 个物体。

  由抽屉原理，至少有两个点落在同一个小线段中。

  而同一小线段的长度不超过 $\dfrac{1}{12}$，所以这两个点之间的距离不超过
  \[
  \frac{1}{12}.
  \]
  \end{solution}
\end{enumerate}

\end{document}
